% Slide "Software prefetch: where the code generator puts it".
% Left: where the pass sits (submodules/pystencils/src/pystencils/codegen/driver.py:313-315: after the
% vector loops are prepared, before the RVV intrinsics are chosen); the pass is
% backend/transformations/insert_prefetches.py, the option cpu.vectorize.prefetch_distance
% (codegen/config.py:278), set by the app (apps/benchmarks/SweepgenLdcCPU/CMakeLists.txt:58,
% --prefetch-distance=128; SweepgenLdcCPU.py:76-77).
% Right, top: the generated StreamCollide of SweepgenLdcCPU_explicit_cse_pf, GCC f32 (board file
% build-gcc-f32/apps/benchmarks/SweepgenLdcCPU/_gen/SweepgenLdcCPU_explicit_cse_pf/gen/SweepgenLdcCPU.cpp,
% lines 96-150), shortened: namespaces dropped, 19 + 23 calls shown as 2 + 1. Right, bottom: the
% runtime function (include/pystencils_runtime/rvv.hpp:144-151, the branch without -march=..._zicbop).
\begin{tikzpicture}
  \tikzset{st/.style={box,text width=36mm,font=\scriptsize,minimum height=5mm,inner sep=2pt}}
  \node[st,quiet] (b1) at (0,0) {lbmpy: the update rule};
  \node[st,quiet,below=2.5mm of b1] (b2) {pystencils: loops, vectorized};
  \node[st,added,below=2.5mm of b2] (b3) {\textbf{InsertPrefetches} (new pass)\\[-1pt]
    {\tiny one prefetch per address, 128 B ahead}};
  \node[st,quiet,below=2.5mm of b3] (b4) {RVV intrinsics, C++ source};
  \node[st,quiet,below=2.5mm of b4] (b5) {compiler: \code{prefetch.r}, \code{.w}};
  \foreach \a/\b in {b1/b2,b2/b3,b3/b4,b4/b5} \draw[arr] (\a) -- (\b);
  \node[codebox,anchor=north west] (gen) at (2.4,0.25) {%
for\ (ctr\_0\ =\ 0;\ ctr\_0\ <\ end;\ ctr\_0\ +=\ 8)\ \{\ \ \textcolor{psgray}{//\ 8\ cells\ per\ step\ (f32)}\\
\ \ prefetch\_read(\&f\_row[ctr\_0\ +\ 32]);\ \ \ \ \ \ \ \ \textcolor{psgray}{//\ 32\ floats\ =\ 128\ B\ ahead}\\
\ \ prefetch\_read(\&f\_row\_\_0[ctr\_0\ +\ 32]);\ \ \ \ \ \textcolor{psgray}{//\ ...\ 19\ read\ prefetches}\\
\ \ prefetch\_write(\&f\_tmp\_row[ctr\_0\ +\ 32]);\ \ \ \textcolor{psgray}{//\ ...\ 23\ write\ prefetches}\\
\ \ xi\_5\ =\ vfadd(vle32(\&f\_row[ctr\_0]),\ ...);\ \ \textcolor{psgray}{//\ then\ the\ update\ ...}\\
\}};
  \node[font=\scriptsize\bfseries,anchor=south west] at (gen.north west) {what the generator writes (names shortened)};
  \node[codebox,anchor=north west] (rt) at ([yshift=-6mm]gen.south west) {%
inline\ void\ prefetch\_read(const\ T\ *p)\ \{\\
\ \ asm\ volatile("ori\ x0,\ \%0,\ 1"\ :\ :\ "r"(p));\ \textcolor{psgray}{//\ =\ prefetch.r}\\
\}};
  \node[font=\scriptsize\bfseries,anchor=south west] at (rt.north west) {the runtime function behind it};
\end{tikzpicture}
